\section{Introduction: The Meta-Question}
\label{sec:intro}

What do we learn about representation design by implementing six alternative
systems and comparing them algorithmically?  This is the meta-question that
animates the E-track research program and that this paper addresses as its
capstone.

The six papers in Track~E --- multi-member districts \citep{dellaLibera2026e1},
county representation \citep{dellaLibera2026e2}, national redistricting
\citep{dellaLibera2026e3}, partisan-similarity clustering
\citep{dellaLibera2026e4}, partisan-fairness allocation \citep{dellaLibera2026e5},
and international applications \citep{dellaLibera2026e6} --- each ask a
local question: what happens to compactness, proportionality, minority
representation, or county preservation when we change this one design dimension?
The synthesis paper E.0 \citep{dellaLibera2026e0} organizes these local
questions into a comparative Pareto analysis.  This paper asks the deeper
question: what general principles about representation system design do we
extract from the systematic comparison?

Our framing is borrowed from engineering design theory, where the concept of
the \emph{design space} refers to the full set of possible configurations of a
system, and the \emph{Pareto frontier} refers to the set of configurations
from which no improvement on one dimension is possible without degrading
another \citep{simon1996sciences}.  We argue that redistricting is a design
problem in exactly this sense: there is a well-defined space of possible
representation systems, a small number of key performance dimensions, and a
set of constraint surfaces that bound the achievable region.  The E-track
experiments are controlled explorations of this space.

\subsection{Why Algorithmic Comparison Enables Design Learning}

The conventional approach to comparing representation systems is political
science field research: examine real-world systems (US single-member districts,
German mixed-member proportional, Irish STV, etc.) and compare outcomes.
This approach has fundamental limitations.  Real-world systems differ on
multiple dimensions simultaneously --- electoral formula, district size, ballot
structure, candidate selection, campaign finance --- making it impossible to
isolate the effect of any single design choice.  They also differ in their
historical, institutional, and demographic contexts, introducing confounds
that threaten causal inference \citep{lijphart1999patterns}.

The algorithmic approach overcomes these limitations by holding the algorithm
constant and varying design dimensions one at a time.  We use the same
edge-weighted recursive bisection procedure \citep{karypis1998fast,
dellaLibera2026mec} for all six experiments.  The algorithm accepts the same
census geography inputs in all cases; what changes is the objective function
(single-member versus multi-member), the jurisdictional constraints (state-based
versus national), the edge-weighting scheme (geographic versus partisan), and
the seat-allocation mechanism (proportional versus winner-take-all).  This
controlled variation means that differences in outcomes can be attributed to
design choices rather than to confounding contextual factors.

A second advantage is counterfactual reachability.  Some alternatives ---
national redistricting, strict county representation --- are constitutionally
unavailable in the US but can be implemented computationally.  The resulting
maps are not real redistricting proposals; they are measurement instruments
that answer the question ``how much does the constitutional constraint cost us
in this dimension?''  Manual redistricting cannot answer this question because
it cannot produce the counterfactual; algorithmic redistricting can.

\subsection{The Five Lessons}

We organize our findings into five design lessons, each addressing a specific
dimension of the design space:

\begin{description}
  \item[Lesson~1 (Section~\ref{sec:lesson1})] \textbf{Geography is destiny for
       single-member districts.}  The Rodden sorting gap \citep{rodden2019why}
       is not a drawing artifact; it is a structural consequence of where
       Democratic voters live.  No single-member geographic drawing method
       eliminates it.

  \item[Lesson~2 (Section~\ref{sec:lesson2})] \textbf{Multi-member districts
       partially but not fully escape geographic sorting.}  Larger districts
       average over more geographic heterogeneity, reducing the wasted-vote
       problem by 40\% (three-member) to 65\% (five-member), but the residual
       gap from extreme urban concentration persists.

  \item[Lesson~3 (Section~\ref{sec:lesson3})] \textbf{Eliminating districts
       entirely trades proportionality for local accountability.}  County
       representation is more proportional in competitive states but less
       responsive to population change and institutionally unfamiliar as a
       voting mechanism.

  \item[Lesson~4 (Section~\ref{sec:lesson4})] \textbf{Crossing state lines
       produces geographically compact but politically fragmented districts.}
       The national map gains 12\% compactness but loses 40\% of whole-county
       preservation and eliminates state-based representation entirely.

  \item[Lesson~5 (Section~\ref{sec:lesson5})] \textbf{The compactness--%
       proportionality Pareto frontier is steep.}  Gaining 1~pp of
       proportionality costs approximately 0.015 PP units.  No system escapes
       this trade-off within a geographic single-election framework.
\end{description}

\subsection{Scope and Method}

All quantitative results in this paper are derived from the redistricting
pipeline described in \citet{dellaLibera2026mec}, applied to 2020 Census
geography for all 50 states and 435 congressional districts.  Political data
(vote shares, efficiency gap, partisan bias) are computed from the 2016--2022
average of US House election results, following the methodology in
\citet{stephanopoulos2015partisan}.  County preservation rates are computed
from Census TIGER 2020 boundary files.  Where results differ from the
companion papers E.1--E.6 (due to data vintage or parameter differences),
we note this explicitly.

The paper targets the audience of \emph{Annual Review of Political Science}:
scholars familiar with redistricting debates and electoral system theory who
seek a synthesizing analysis rather than the primary experimental findings.
Readers seeking methodological detail should consult the companion papers.
